\section{Rigid inner forms and $\Bun_{G}^e$}\label{sec:rigid-inner-forms}


In this section we recall some fundamental geometric objects used in the paper.

\subsection{The gerbe $\Kal$}
We first recall the reinterpretation of $\Kal$ from \cite{Fargues22} as a gerbe over the Fargues--Fontaine curve.

\begin{defi}
\begin{enumerate}
\item{Define the pro-multiplicative group scheme $\mathbb{D} := \varprojlim_{n} \mathbb{G}_{m}$.}
    \item{We denote by $\Kott_{F}$ the Kottwitz gerbe (this is sometimes denoted by $\mc{E}^{\tn{iso}}$ and called the ``isocrystal gerbe'' in the literature but we choose $\Kott_{F}$ in order to be consistent with \cite{Fargues22}), which is the gerbe over $\Spec(F)$ corresponding to the Tannakian category $\Isoc_{F}$ of $F$-isocrystals, which is banded by $\mathbb{D}$.}

\end{enumerate}
    We will frequently just write $\mathrm{Kott}$ for $\mathrm{Kott}_{F}$ and reserve the notation $\mathrm{Kott}_{E}$ for the analogous object for a finite Galois extension $E/F$. We also write $\mathrm{Kott}$ to denote the corresponding \'{e}tale gerbe over $\mathrm{Spa}(F)$.
\end{defi}

\begin{lem}\label{lem:maptoKott}
    There is a canonical morphism over $\Spa(F)$
    \begin{equation*}
        X_{S} \to \Kott
    \end{equation*}
    for any $S \in \Perfd$.
\end{lem}

\begin{proof}
For any $S \in \Perfd$ there is an exact, faithful tensor functor
\begin{equation*}
\tn{Isoc}_{F} \to \tn{VB}(X_{S}),
\end{equation*}
where $\tn{VB}(X_{S})$ denotes the tensor category of vector bundles on $X_{S}$ (by \cite[\S III.2.1]{Geometrization}). We are then done because $\Kott$ is the gerbe associated to the Tannakian category $\Isoc_{F}$. 
\end{proof}

Let $E/F$ be a finite Galois extension. As explained in \cite[\S 4.1]{Fargues22}, we can define an isocrystal of slope $-1$ over $E$ as follows. First, one chooses an arbitrary slope-one isocrystal $(D,\varphi)$ over $E$, and then defines the $\breve{E}$-vector space
\begin{equation*}
V(D,\varphi) = \{f \colon D \to B_{\tn{cris}}^{(E)} | f \circ \tn{Frob}_{E} = \tn{Frob}_{E} \circ f, \hspace{1mm} f(D) \subseteq \tn{Fil}^{1} B_{\tn{cris}}^{(E)}\}
\end{equation*}
where $B_{\tn{cris}}^{(E)}$ is as defined in \cite[\S 1.2]{FF} (along with its natural filtration and Frobenius action). The desired isocrystal is defined as
\begin{equation*}
    \mathbb{L}_{E} := (D,\varphi)^{\vee} \otimes_{\breve{E}} V(D,\varphi)^{\vee}. 
\end{equation*}
The key property of $\mathbb{L}_{E}$ for our purposes is that if one replaces $(D,\varphi)$ with another slope-one isocrystal which produces $\mathbb{L}_{E}'$ via the above recipe, then there is a canonical isomorphism from $\mathbb{L}_{E}$ to $\mathbb{L}_{E}'$ (cf. \cite[\S 4.1]{Fargues22} for more details). Further, it is shown in Lemme 4.3 loc. cit. that, for $E'/E/F$ two such extensions, there is a canonical identification between 
\begin{equation*}
    N_{E'/E} \mathbb{L}_{E'} := \bigotimes_{\sigma \in \Gal_{E'/E}} \sigma^{*}(\mathbb{L}_{E'}) \cong \mathbb{L}_{E}.
\end{equation*}


\begin{defi}\label{def:LE/F}
The datum of $\mathbb{L}_{E}$ defines a $\mathbb{G}_{m}$-torsor on $\mathrm{Kott}_{E}$. Recalling that $X_{S}^{(E)} \to X_{S}$ denotes the corresponding finite \'{e}tale cover of $X_{S}$, we can pull this torsor back via the map from Lemma \ref{lem:maptoKott} to a $\mathbb{G}_{m}$-torsor over $X_{S}^{(E)}$, denoted by $\mc{T}_{E/F,S}$. We can also pull the line bundle $\mathbb{L}_{E}$ back to $X_{S}$, which we denote by $\mc{L}_{E/F,S}$. Note that $\mc{T}_{E/F,S}$ is the sheaf of trivializations of $\mc{L}_{E/F,S}$. 
\end{defi}

We define (as in \cite[\S 3.1]{Kal16}, where we are using the version without quotients, see \cite[Remark 2.12]{BMD26} for a detailed explanation of the difference)
\begin{equation*}
u = \varprojlim_{E/F,n} \Res_{E/F}(\mu_{n}),
\end{equation*}
where the transition maps are given by field norms and $n/m$-power maps (for $m \mid n$). Similarly, set $\tilde{t} = \varprojlim_{E/F,n} \Res_{E/F}(\mathbb{G}_{m})$ and $t = \varprojlim_{E/F} \Res_{E/F}(\mathbb{G}_{m})$, which fit into the short exact sequence of multiplicative group schemes over $F$
\begin{equation*}
1 \to u \to \tilde{t} \to t \to 1.
\end{equation*}

The sequence of $\mathbb{G}_{m}$-torsors on each $\mathrm{Kott}_{E}$ obtained from $\{\mathbb{L}_{E}\}_{E/F}$ defines a $t$-torsor $\mathbb{T}$ on $\mathrm{Kott}_{F} =: \mathrm{Kott}$, and we can form the $u$-gerbe 
\begin{equation*}
[\mathbb{T}/\tilde{t}] \to \mathrm{Kott},
\end{equation*}
which is also a gerbe over $\Spec(F)$.

\begin{rque}
    The above construction of the system $\{\mathbb{L}_{E}\}_{E/F}$ is a nontrivial part of \cite{Fargues22} relying on ideas from $p$-adic Hodge theory, and we do not recall any details beyond the above summary and refer the reader to \cite[\S 4.1]{Fargues22} for a detailed discussion. The only properties of the system $\{\mathbb{L}_{E}\}_{E/F}$ that we use here are that it consists of isocrystals of slope $-1$, is unique up to canonical isomorphism, and is norm-compatible. In fact, for the purpose of understanding this paper, it is essentially harmless to replace the system $\{\mathbb{L}_{E}\}_{E/F}$ with the system $\{(D_{E}, \varphi_{E})\}$ of isocrystals determined by picking a norm-compatible system of uniformizers $\{\varpi_{E}\}_{E/F}$ for all finite Galois $E/F$. Using this system instead will produce the desired spectral action---the subtle point is that the resulting action will only be canonical up to acting on $\Bun_{G}^{e}$ (cf. Definition \ref{defi:BunGe}) via pulling back by the automorphism of $[\mathbb{T}/\tilde{t}]$ determined by the differential of a $0$-cochain valued in $u(\ov{F})$.
    
    Furthermore, since in \S \ref{sec:extendedheckedef} we need to choose a system of norm-compatible untilts $\{y^{(E)}\}_{E/F}$ anyway, one could just as well take the system of isocrystals to be the one obtained by taking the fiber of each $\mc{O}_{X_{S}^{(E)}}(y^{(E)})$ at the closed point of $X_{S}^{(E)}$ corresponding to $y^{(E)}$. In this way, one obtains a norm-compatible system of isocrystals of slope $-1$ without changing the number of choices made in this paper. 
    \end{rque}

\begin{defi}\label{defi:Tdef}
We denote by $\mc{T}_{S}$ the pro-\'{e}tale $t_{S}$-torsor over $X_{S}$ corresponding to the norm-compatible system of $\mathbb{G}_{m}$-torsors $\{\mc{L}_{E/F,S}\}_{E/F}$.
\end{defi}

We then obtain a $u$-gerbe over $X_{S}$ by taking the quotient stack
\begin{equation*}
    [\mc{T}_{S}/\tilde{t}_{S}] \to X_{S}.
\end{equation*}
When $F$ has mixed characteristic, this is a pro-\'{e}tale gerbe. When $F$ has equal characteristic, it is a gerbe for the v-topology. By \cite[Corollaire 7.4]{Fargues22}, there is a canonical identification
\begin{equation*}
    [\mc{T}_{S}/\tilde{t}_{S}] \xrightarrow{\sim} X_{S} \times_{\Kott} [\mathbb{T}/\tilde{t}].
\end{equation*}

We conclude this subsection by explaining how the gerbes $[\mc{T}_{S}/\tilde{t}_{S}]$ and $[\mathbb{T}/\tilde{t}]$ are related to Kaletha's original construction in \cite{Kal16}. 

\begin{prop}\label{prop:ucohom}
   We have the following canonical isomorphisms:
    \begin{enumerate}
       \item $H^{1}_{\tn{fppf}}(F, u) = 0$;
     \item $H^{2}_{\tn{fppf}}(F,u) = \widehat{\Z}$.
   \end{enumerate}
\end{prop}

\begin{proof}
    This is \cite[Theorem 3.1]{Kal16} in the mixed characteristic case and \cite[Theorem 3.4]{Dillery23} in the equal characteristic case.
\end{proof}

The Kaletha class is defined as the class $[\Kal] \in H^{2}_{\tn{fppf}}(F, u)$ corresponding to $-1$ under the above identification. 

\begin{prop}
\begin{enumerate}
\item{The stack 
\begin{equation*}
[\mathbb{T}/\tilde{t}] \to \mathrm{Kott} \to \Spec(F)
\end{equation*}
is a $\mathbb{D} \times u$-gerbe and represents the class $[\mathrm{Kott}] \times [\Kal] \in H_{\tn{fppf}}^{2}(F, \mathbb{D} \times u)$. In particular, the $u$-gerbe 
\begin{equation*}
[\mathbb{T}/\tilde{t}] \times^{B(\mathbb{D} \times u)} Bu\to \Spec(F)
\end{equation*}
represents $[\Kal]$; denote this gerbe by $\Kal$.}

    \item{There is a canonical isomorphism of stacks over $X_{S}$
    \begin{equation*}
        [\mc{T}_{S}/\tilde{t}_{S}] \xrightarrow{\sim} X_{S} \times_{\tn{Spa}(F)} \tn{Kal}.
    \end{equation*}}
    \end{enumerate}
\end{prop}

\begin{proof}
    This follows from \cite[Th\'{e}or\`{e}me 5.6]{Fargues22} and \cite[Corollaire 7.4]{Fargues22}.
\end{proof}

\begin{rque}
    One can also directly define a $\tilde{t}$-gerbe $\mc{E}$ over $\Spec(F)$ by considering the direct limit (over all $E/F$ finite Galois and $n \in \mathbb{N}$) of the Tannakian categories $\mc{I}_{E/F,n}$ defined in \cite[\S 2.5]{BMD26}. It has the property that $\mc{E} = \Kal \times^{Bu} B\tilde{t}$.
\end{rque}

\subsection{The stack $\tn{Bun}_{G}^{e}$}
Let $G$ be an affine group scheme over $F$; we denote by $G_{S}$ the corresponding sheaf of groups on $X_{S}$ for $S \in \Perfd$. We equip the gerbe $[\mc{T}_{S}/\tilde{t}_{S}] \to X_{S}$, which is a v-stack over $S$, with its natural étale topology, and denote by $G_{S}$ the pullback of $G_{S}$ to a group sheaf on  $[\mc{T}_{S}/\tilde{t}_{S}]$ for this topology. When we talk about $G_{S}$-torsors on $[\mc{T}_{S}/\tilde{t}_{S}]$, we always mean with this \'{e}tale topology. 

\begin{rque} Occasionally we will talk about $Z_{S}$-torsors over $[\mc{T}_{S}/\tilde{t}_{S}]$ for $Z$ a finite central subgroup of $G$, and in such cases will use the v-topology on $[\mc{T}_{S}/\tilde{t}_{S}]$ when $F$ has equal characteristic. This situation will only arise when a $Z_{S}$-torsor for the v-topology acts on an \'{e}tale $G_{S}$-torsor by taking contracted products, which yields another \'{e}tale $G_{S}$-torsor and is thus completely harmless for our purposes.  
\end{rque}

Given a $G_{S}$-torsor $\cP$ on  $[\mc{T}_{S}/\tilde{t}_{S}]$, the action of the band determines a morphism of sheaves over $[\mc{T}_{S}/\tilde{t}_{S}]$
\begin{equation}\label{eq:bandmap}
    u_{S} \to \underline{\tn{Aut}}_{[\mc{T}_{S}/\tilde{t}_{S}]}(\cP)/G_{\tn{ad},S},
\end{equation}
see \cite[Lemma 2.27]{Dillery23} for more details. 

\begin{defi}\label{defi:bandmap}
    We say that a $G_{S}$-torsor $\cP$ on $[\mc{T}_{S}/\tilde{t}_{S}]$ is \textbf{Kal-basic} if the homomorphism \eqref{eq:bandmap} is obtained by composing a (uniquely-determined) homomorphism 
    \begin{equation*}
    \lambda_{\cP} \colon u_{S} \to Z_{G,S}
    \end{equation*}
    with the action of $Z_{G,S}$ on $\cP$ induced by its structure as a $G_{S}$-torsor.

    A priori, the map $\lambda_{\cP}$ is a morphism of group schemes over $[\mc{T}_{S}/\tilde{t}_{S}]$, but by \cite[Lemma 2.42]{Dillery23}, this canonically descends to a morphism of group schemes over $X_{S}$, which we also denote by $\lambda_{\cP}$.
\end{defi}

For $S \in \Perfd$ and a $G$-torsor $\cP \in \Bun_{G}^{e}(S)$ we obtain, by construction, a morphism $\lambda_{\cP} \in \Hom_{X_{S}}(u_{S}, Z_{G,S})$, which we call the \textbf{inertial homomorphism} of $\cP$. 

\begin{defi}\label{defi:BunGe}
\begin{enumerate}
   \item{For $G$ an affine group scheme over $F$ (usually taken to be a connected reductive group), we define $\Bun_{G}^{e}$ to be the functor on $\Perfd$ sending $S$ to the groupoid of all $\Kal$-basic $G_{S}$-torsors on $[\mc{T}_{S}/\tilde{t}_{S}]$.}

\item{For $Z$ a finite central subgroup of $G$, we define $\Bun_{Z \subset G}^{e}$ to be the subfunctor of $\Bun_{G}^{e}$ of $G$-torsors such that $\lambda_{\cP}$ factors through $Z_{S}$ for $\cP \in \Bun_{G}^{e}(S)$.}
   \end{enumerate}
\end{defi}

\begin{rque}
   It is clear that $\Bun_{G}^{e} = \varinjlim_{Z} \Bun_{Z \subset G}^{e}$. 
\end{rque}

\begin{defi}[\protect{\cite[D\'{e}finition 9.1]{Fargues22}}]\label{defi:BeGdef}
    For any affine group scheme $G$ over $F$, define the extended Kottwitz set, denoted by $B_{e}(G)$, to be the pointed set of isomorphism classes of all $\Kal$-basic $G$-torsors on $\mathrm{Kott} \times_{F} \Kal$ (where $\Kal$-basic has the analogous meaning as above). For finite central $Z$ in $G$ we define $B_{e}(Z \subset G)$ analogously. Define $B_{e}(G)_{\tn{basic}}$ to be the subset of torsors whose restriction to $\mathbb{D} \times u$ factors through $Z_{G}$.
\end{defi}

We now summarize a few key results about $\Bun_{G}^{e}$:

\begin{prop}[{\cite[Proposition 12.5, Exemple 12.6]{Fargues22}}]\label{prop:BunGestructure}
Let $G$ be a connected reductive group over $F$.
    \begin{enumerate}
       \item{The functors  $\Bun_{G}^{e}$ and  $\Bun_{Z \subset G}^{e}$ are small $v$-stacks.}
       \item{There are canonical identifications
       \begin{equation*}
           |\Bun_{G}^{e}| \xrightarrow{\sim} B_{e}(G)
       \end{equation*}
       and
       \begin{equation*}
            |\Bun_{Z \subset G}^{e}| \xrightarrow{\sim} B_{e}(Z \subset G).
       \end{equation*}       
       }
       \item{For a fixed $\lambda \in \Hom_{F}(u, Z_{G})$, define the substack $\Bun_{G}^{e,\lambda}$ to consist of all $\cP$ such that $\lambda_{\cP_{\bar{s}}} = \lambda$ for all geometric points $\bar{s}$ of $S$. We have a decomposition (as stacks)
       \begin{equation*}
           \Bun_{G}^{e} = \bigsqcup_{\lambda \in \Hom_{F}(u,Z_{G})} \Bun_{G}^{e,\lambda}.
       \end{equation*}
       The same decomposition holds for $\Bun_{Z \subset G}^{e}$ if we replace $Z_{G}$ by $Z$.

       Moreover, if $b \in B_{e}(G)$ is an element whose restriction to $u$ is $\lambda$ (which always exists), then twisting by $b$ gives an isomorphism of stacks
       \begin{equation*}
           \Bun_{G}^{e,\lambda} \xrightarrow{\sim} \Bun_{G_{b}}.
       \end{equation*}
       }
\item{For a fixed $b \in B_{e}(G)_{\tn{basic}}$, define the substack $\Bun_{G}^{e,b}$ to consist of all $\cP$ such that $\cP_{\bar{s}} \cong b$ (using $\tn{(2)}$) for all geometric points $\bar{s}$ of $S$. Then the semistable locus of $\Bun_{G}^{e}$ equals (as a stack) $\bigsqcup_{b \in B_{e}(G)_{\tn{basic}}}\Bun_{G}^{e,b}$.}

    \end{enumerate}
\end{prop}

\begin{proof}
    The first two statements are \cite[Proposition 12.5]{Fargues22} and the third and fourth follow immediately from Exemple 12.6 loc. cit.
\end{proof}

\begin{lem}\label{lem:inertialhom}
    We have a canonical isomorphism
    \begin{equation*}
        \Hom_{X_{S}}(u_{S}, Z_{G,S}) \xrightarrow{\sim} \prod_{\pi_{0}(|X_{S}|)} \Hom_{F}(u,Z_{G}).
    \end{equation*}
\end{lem}

\begin{proof}
First, we observe that any morphism 
 \begin{equation*}
u_{S} \times_{\Spa(F)} X_{S} \to Z_{G,S} \times_{\Spa(F)} X_{S}
 \end{equation*}
 is locally constant over $|X_{S}|$, as in the proof of \cite[Proposition 12.5]{Fargues22}, which allows us to assume that $|X_{S}|$ is connected. The same proof shows that any such morphism is locally constant over $|S|$, which reduces us further to the case $S = \Spa(C,C^{+})$, where it follows from the fact that the étale fundamental group of $X_{C}$ is $\Gal_{F}$, by \cite[Th\'{e}or\`{e}me 9.5.1]{FF}.
\end{proof}

\begin{rque}\label{rque:fingerbe}
\begin{enumerate} \item{Observe that, for any $S \in \Perfd_{C}$, the fact that $G$ is finite type implies that any $G_{S}$-torsor on $[\mc{T}_{S}/\tilde{t}_{S}]$, $[\mc{T}_{S}^{\circ}/\tilde{t}_{S}]$, $[\mc{T}_{S}^{\loc}/\tilde{t}_{S}]$, or $[\mc{T}_{S}^{\loc,\circ}/\tilde{t}_{S}]$ is pulled back from a finite-level version of $[\mc{T}_{S}/\tilde{t}_{S}]$ (respectively its punctured and/or local analogues) obtained by replacing $\mc{T}_{S}$ by $\mc{T}_{E/F,S}$ (as in Definition \ref{def:LE/F}) and $\tilde{t}_{S}$ by $\tilde{t}_{E/F,n,S}$, where $\tilde{t}_{E/F,n} = \Res_{E/F}(\mathbb{G}_{m}) \to  \Res_{E/F}(\mathbb{G}_{m}) = t_{E/F}$ is the covering by the $n$-power map with kernel $u_{E/F,n}$.}

\item{It will be useful for some calculations to choose, for once and for all, an $\ov{F}$-trivialization of the gerbe  $[\mathbb{T}/\tilde{t}] \to \Kott$, which determines a \v{C}ech $2$-cocycle $\xi \in u(\ov{F}^{\bigotimes_{F}3})$ and which we can pull back to a $2$-cocycle for the gerbe $[\mc{T}_{S}/\tilde{t}_{S}]$, also denoted by $\xi$. This determines analogous $2$-cocycles $\xi_{E/F,n}$ for each finite level $u_{E/F,n}$-gerbe $[\mc{T}_{E/F,n,S}/\tilde{t}_{E/F,n,S}]$ in (1)---this cocycle is a priori valued in $u_{E/F,n}(\ov{F}^{\bigotimes_{F}3})$, but in fact takes values in $u_{E/F,n}(E')$ for finite Galois $E'/F$ trivializing $[\mc{T}_{E/F,n,S}/\tilde{t}_{E/F,n,S}]$ (for example, choose $E \subset E'$ such that $n \mid [E' \colon E]$).}

\end{enumerate}
\end{rque}

Since $\Bun_{G}^{e}$ is a v-stack, it makes sense to define, for $\Lambda$ a $\Z_{\ell}$-algebra, the category $\mc{D}_{\tn{lis}}(\Bun_{G}^{e}, \Lambda)$ as in \cite[Definition VII.6.1]{Geometrization} in this context.

\begin{corol}\label{cor:Dlise}
\begin{enumerate}
\item{For $b \in B_{e}(G)$, we have a canonical morphism of v-stacks $\Bun_{G}^{e,b} \to [\pt/G_{b}(F)]$. Moreover, the pullback functors
\begin{equation*}
    \mc{D}_{\tn{lis}}(\Bun_{G}^{e,b}, \Lambda) \to \mc{D}_{\tn{lis}}([\pt/G_{b}(F)], \Lambda) \to \mc{D}_{\tn{lis}}([\Spa(C)/G_{b}(F)], \Lambda)
\end{equation*}
are equivalences, and all of these categories are naturally equivalent to the category of smooth representations of $G_{b}(F)$ on discrete $\Lambda$-modules.}
\item{Pullback induces an equivalence
\begin{equation*}
\mc{D}_{\tn{lis}}(\Bun_{G}^{e}, \Lambda) \to \mc{D}_{\tn{lis}}(\Bun_{G}^{e} \times \Spd(C), \Lambda),
\end{equation*}
similarly with $\Bun_{G}^{e}$ replaced by $\Bun_{G}^{e,\lambda}$ for $\lambda \in \Hom_{F}(u,Z)$ or $\Bun_{G}^{e,b}$ for $b \in B_{e}(G)_{\tn{basic}}$.
}
\end{enumerate}
\end{corol}

\begin{proof}
    This result comes from combining Proposition \ref{prop:BunGestructure} with \cite[Proposition VII.7.1]{Geometrization} and \cite[Proposition VIII.7.3]{Geometrization}.
\end{proof}

\subsection{Connected components}
Recall from \cite[D\'{e}finition 9.10]{Fargues22} that one has an ``extended'' version of the Galois coinvariants of the fundamental group, defined as follows. Recall that $\Phi$ denotes the (absolute) root system of $T$ (with coroots $\Phi^{\vee}$) and set $\langle \Phi \rangle^{*} := \{\nu \in X_{*}(T)_{\mathbb{Q}} | \langle \alpha, \nu \rangle \in \Z \hspace{1mm} \forall \alpha \in \Phi\}$. The group in question is given by
\begin{equation*}
\pi_{1}(G)^{e}_{\Gal_{F}} := \frac{\langle \Phi \rangle^{*}}{\langle \Phi^{\vee} \rangle + I_{\Gal_{F}} \cdot X_{*}(T)},
\end{equation*}
where $I_{\Gal_{F}}$ denotes the augmentation ideal for $\Gal_{F}$.

Recall from \cite[\S 9.5]{Fargues22} that we have an exact sequence
\begin{equation}\label{eq:pi1eSES}
    0 \to \pi_{1}(G)_{\Gal_{F}} \to \pi_{1}(G)^{e}_{\Gal_{F}}  \to \Hom_{F}(u,Z_{G}) \to 0,
\end{equation}
and also that there is an ``extended Kottwitz map'' (recalling Definition \ref{defi:BeGdef})
\begin{equation*}
    B_{e}(G) \xrightarrow{\kappa}  \pi_{1}(G)^{e}_{\Gal_{F}}
\end{equation*}
whose restriction to $B_{e}(G)_{\tn{basic}}$ is a bijection (thus giving the left-hand side the structure of an abelian group) and whose restriction to $B(G)$ recovers the usual Kottwitz map (see \cite[\S 9.5.3]{Fargues22} for these claims).

It follows from Proposition \ref{prop:BunGestructure} that there is a map
\begin{equation*}
    |\Bun_{G}^{e}| \xrightarrow{\kappa} \pi_{1}(G)^{e}_{\Gal_{F}}
\end{equation*}
which Fargues proves (\cite[Proposition 12.5]{Fargues22}) is locally constant. 

By combining Proposition \ref{prop:BunGestructure} with the proof of \cite[Proposition VII.7.3]{Geometrization} and grouping together the fibers of $\kappa$ by their image along the surjection $\pi_{1}(G)^{e}_{\Gal_{F}}  \to \Hom_{F}(u,Z_{G})$, we obtain:
\begin{prop} \begin{enumerate}
   \item{There is an infinite semi-orthogonal decomposition of  $\mc{D}_{\tn{lis}}(\Bun_{G}^{e},\Lambda)$ into $\mc{D}_{\tn{lis}}(\Bun_{G}^{e,b},\Lambda)$ as $b$ ranges over all $B_{e}(G)$. In particular, we have  
    \begin{equation*}
        \mc{D}_{\tn{lis}}(\Bun_{G}^{e},\Lambda) \xrightarrow{\sim} \prod_{x \in \pi_{1}(G)^{e}_{\Gal_{F}}}  \mc{D}_{\tn{lis}}(\Bun_{G}^{e,\kappa = x},\Lambda).
    \end{equation*}}
    \item{For a fixed $\lambda \in  \Hom_{F}(u,Z_{G})$, there is a semi-orthogonal decomposition \begin{equation*}
\mc{D}_{\tn{lis}}(\Bun_{G}^{e,\lambda},\Lambda) =  \prod_{x \in \pi_{1}(G)^{e}_{\Gal_{F}}, x|_{u} = \lambda}  \mc{D}_{\tn{lis}}(\Bun_{G}^{e,\kappa = x},\Lambda) \xrightarrow{\sim} 
 \prod_{x \in \pi_{1}(G_{b_{\lambda}})_{\Gal_{F}}}  \mc{D}_{\tn{lis}}(\Bun_{G_{b_{\lambda}}}^{\kappa = x},\Lambda),
\end{equation*}
   where $b_{\lambda} \in B_{e}(G)$ lifts $\lambda$ and the second isomorphism is obtained via twisting (cf. Proposition \ref{prop:BunGestructure}).}
   \end{enumerate}
\end{prop}

\subsection{An affine Grassmannian calculation}\label{sec:affGrasscalc}
The purpose of this subsection is to provide some calculations which will be used in \S \ref{sec:extendedheckedef}.

Consider the affine Grassmannian 
\begin{equation*}
\mathrm{Gr}_{t} := \varinjlim_{E/F} \tn{Gr}_{\tn{Res}_{E/F}(\mathbb{G}_{m})}
\end{equation*}
for $t$, as defined in \cite[\S 19]{SW20} (this direct limit is harmless because sheafification commutes with filtered colimits). A compatible system of untilts $\{\Spa(C^{\sharp},C^{\sharp,+}) \xrightarrow{{y^{(E)}}} X_{C}^{(E)}\}_{E/F}$, which we have fixed (cf. \S \ref{sec:choiceoflifts} for a discussion of this choice), provides an identification


\begin{equation*}
    \mathrm{Gr}_{t}(C) \xrightarrow{\sim} X_{*}(t),
\end{equation*}
and we consider the homomorphism of abelian groups $X_{*}(t) \ni \tn{Sh}\colon X^{*}(t) = \mathcal{C}(\Gal_{F},\Z) \to \Z$ sending a function $f$ to $f(e)$ (here $\tn{Sh}$ stands for ``Shapiro'', which we use because the isomorphism $\Hom_{F}(t,\mathbf{T}) \xrightarrow{\sim} X_{*}(\mathbf{T})$ sends $f$ to $f \circ \tn{Sh}$). 

Denote by $\bbD_{C}$ the formal disk $\Spec(B_{\tn{dR}}^{+}(C^{\sharp}))$ corresponding to $y$ and $\bbD_{C}^{\circ}$ the associated punctured disk. 

\begin{lem}\label{lem:Grt} There is a unique trivialization $\psi \colon \mc{T}_{C}^{\circ} \xrightarrow{\sim} \underline{t_{C}}$ whose image in $\tn{Gr}_{t}(C)$ (by pulling back to $\bbD_{C}$) is $\tn{Sh}$. 
\end{lem}

\begin{proof}
    For a fixed finite Galois $E/F$, the torsor $\mc{T}_{E,C}$ is the sheaf of trivializations of $\mathcal{L}_{E/F,C}$ (as in Definition \ref{def:LE/F}), viewed as a sheaf of $\mathcal{O}_{X_{C}}$-modules. The system $\{y^{(E)}\}_{E/F}$ of untilts determines an isomorphism of projective systems of line bundles 
    \begin{equation*}
    \{\mc{L}_{E/F,C}\}_{E/F} \xrightarrow{\sim} \{\mc{O}_{X^{(E)}_{C}}(y^{(E)})\}.
    \end{equation*}
    
    Therefore, to trivialize $\mc{T}_{E,C}$ it suffices to find a norm-compatible system of trivializations of each line bundle $\mc{O}_{X^{(E)}_{C}}(y^{(E)})$ and, in turn, a norm-compatible system of uniformizing parameters $\xi_{E}$ for each $y^{(E)}$. Note that each $\xi_{E}$ can be chosen to lie in $\mathcal{O}_{X_{C}^{(E)}}(y^{(E)})(X^{(E)}_{C} \setminus \Gamma_{y^{(E)}})$, and so the resulting trivialization $\psi$ of $\mc{T}_{C}^{\circ}$ is defined on $X_{C}^{\circ}$. 

    Now $\psi$ has the correct image in $\tn{Gr}_{t}(C)$, because, by definition of the isomorphism 
    \begin{equation*}
    \tn{Gr}_{\Res_{E/F}(\mathbb{G}_{m})}(C)\xrightarrow{\sim} \Hom_{\Z}(\Z[\Gal_{E/F}],\Z),
    \end{equation*}
    the element $\tn{Sh}$ corresponds to the automorphism of the trivial torsor on $\bbD_{C}^{\circ}$ sending $e$ to $\xi_{E}^{-1} \in \Res_{E/F}(\mathbb{G}_{m})(\bbD_{C}^{\circ})$, which is what our map $\psi$ yields (after picking an auxiliary trivialization of $\mc{T}_{C}$ on $\bbD_{C}$ by using some closed point of $X_{C}$ different from the image of $y$). Taking the projective limit over $E/F$ gives the result.

   For uniqueness, we observe that the collection of $X_{C}^{\circ}$-trivializations of $\mc{T}_{C}^{\circ}$ is a $t(X_{C}^{\circ})$-torsor, and our choice of untilts $\{y^{(E)}\}$ determines an isomorphism 
   \begin{equation*}
       t(X_{C}^{\circ}) \xrightarrow{\sim} \varprojlim_{E/F} \Z[\Gal_{E/F}]_{0},
   \end{equation*}
   where the subscript ``$0$'' denotes the kernel of the augmentation map. Note that $X_{*}(t) = \varprojlim \Z[\Gal_{E/F}]$ and, via the identification $\tn{Gr}_{t}(C) \xrightarrow{\sim} X_{*}(t)$ fixed above, replacing $\psi$ with an $x\in t(X_{C}^{\circ})$-translate sends it to $\tn{Sh} \cdot f_{x} \in X_{*}(t)$, where $f_{x} \in X_{*}(t)$ is the cocharacter corresponding to $x$.
\end{proof}

\begin{lem}\label{lem:cohomvanish}
    For any algebraically closed perfectoid field $C'$, the cohomology groups $H^{1}(\bbD_{C'}, t)$ and $H^{1}(\bbD_{C'}^{\circ}, t)$ are trivial. In particular, the torsor $\mc{T}_{C'}$ has trivializations on $\bbD_{C'}$.
\end{lem}

\begin{proof}
    We only give the proof for $H^{1}(\bbD_{C'}^{\circ}, t)$; the other case is trivial by Henselianity.
    
    Note that the map $\Res_{E'/F}(\mathbb{G}_{m})(\bbD_{C'}^{\circ}) \to \Res_{E/F}(\mathbb{G}_{m})(\bbD_{C'}^{\circ})$ is surjective, because the kernel $K$ of this map is a torus and $H^{1}(\bbD_{C'}^{\circ}, K) =0$, due to the fact that $\Bdr(C')$ is a complete discrete valuation field with algebraically closed residue field (by, e.g., \cite[\S III.2.3]{Serre94}). 

    In particular, we have $\varprojlim^{(1)} H^{0}(\bbD_{C'}^{\circ}, \Res_{E/F}(\mathbb{G}_{m})) =0$, and thus we have an isomorphism
    \begin{equation*}
        H^{1}(\bbD_{C'}^{\circ}, t) \xrightarrow{\sim} \varprojlim H^{1}(\bbD_{C'}^{\circ}, \Res_{E/F}(\mathbb{G}_{m})) = 0,
    \end{equation*}
    where the last equality is again from the fact that $\Bdr(C')$ is a complete discrete valuation field with algebraically closed residue field. 
\end{proof}

Fix also $\tilde{\psi}$ an auxiliary trivialization of $\mc{T}_{C}^{\loc}$ on $\bbD_{C}$, which exists by the above lemma.